\DocumentMetadata{
  pdfversion=2.0,
  lang=es-MX,
  pdfstandard=ua-2
}

\documentclass{sener2025}

\addbibresource{referencias.bib}

% --- Metadatos PDF/UA (Accesibilidad Universal) ---
\hypersetup{
  pdftitle={Balance Nacional de Energía},
  pdfauthor={SENER},
  pdfsubject={BNE 2025},
  pdfkeywords={Energía},
  pdfcreationdate={D:20251218002019},
  pdfversion={1}
}

% --- Metadatos del Documento ---
\title{Balance Nacional de Energía}
\subtitle{BNE 2025}
\author{SENER}
\date{17 de diciembre de 2025}
\institucion{Secretaría de Energía (SENER)}
\unidad{Unidad de Planeación y Transición Energética}
\setDocumentoCorto{Balance Nacional de Energía 2024}
\palabrasclave{Energía}
\version{1}

\begin{document}

\portadafondo[img/portada.jpg]

\tableofcontents
\newpage

\listatablas
\newpage

\clearpage
\begin{center}
{\Large\patriafont\bfseries\color{gobmxGuinda}Presentación}\\[1cm]
\end{center}

La Seguridad Energética es un objetivo estratégico del Plan Nacional de Desarrollo 2025-2030, pero debe hacerse con una visión de sustentabilidad ambiental; para ello es fundamental consolidar la rectoría del Estado en el sector energético mediante el fortalecimiento de PEMEX y de la CFE. Esto permitirá eliminar la dependencia del exterior, asegurar precios accesibles para la población y avanzar hacia la autosuficiencia energética. En la matriz energética del país se impulsarán fuentes de energía renovables y se acelerará la transición energética con el objetivo de reducir las emisiones contaminantes, cumplir con las metas nacionales de energías limpias y honrar los compromisos internacionales en la lucha contra el cambio climático.

Además de garantizar un suministro energético seguro, eficiente y asequible para toda la población, la justicia energética debe convertirse en un pilar del desarrollo nacional. Esto implica ampliar la cobertura y el acceso a la energía, especialmente en comunidades marginadas, asegurando que todas las regiones del país cuenten con fuentes de energía limpias y sostenibles.

En México se ha tomado la decisión de trabajar con una visión de preservar la seguridad y autosuficiencia energética, encaminarse a una transición energética justa y sustentable, así como de brindar certeza jurídica a las actividades que se realizan en el sector. Por estas razones, se han hecho cambios al marco jurídico del sector energético partiendo de las modificaciones a la Constitución Política de los Estados Unidos Mexicanos y a diversos ordenamientos jurídicos que de ella emanan. En este sentido, la SENER tiene el mandato legal de ser el órgano rector de la política y planeación energética nacional, bajo criterios de soberanía y seguridad energética, autosuficiencia, restitución y aumento de reservas, diversificación de fuentes de combustibles, mejoramiento de la productividad, la reducción progresiva de impactos ambientales en la producción y consumo de energía, y la satisfacción de las necesidades energéticas básicas mediante cobertura universal a toda la población a precios accesibles.

\portadaseccion{1}{Introducción}{}

\section{Introducción}

El sector energético es un área estratégica para el funcionamiento de la economía y el bienestar de la sociedad mexicana, ya que proporciona energía para el desarrollo nacional a través de diversos sectores productivos (agropecuario, industrial, residencial, comercial, transporte, de servicios públicos, entre otros). También, suministra insumos para fines no energéticos en la industria petroquímica, como es la producción de fertilizantes, medicamentos y plásticos ampliamente utilizados en un sin número de artículos de uso cotidiano.
Con el propósito de comprender y analizar la dinámica energética del país, así como de generar la base para la planeación vinculante, la Secretaría de Energía (SENER) elabora anualmente el Balance Nacional de Energía (BNE). Actualmente, este instrumento se desarrolla en el marco de la Reforma Energética de 2025, lo que le confiere un papel fundamental en la definición de políticas y estrategias para el Sector Energético Mexicano (SEM).
En este documento no solo se analiza el desempeño del SEM en el año 2024, sino que se detalla el periodo 2010 - 2024 a fin de comprender su comportamiento histórico. Por lo tanto, este documento describe las fuentes de energía, su extracción y producción, transformación a productos energéticos secundarios y el consumo final de estos. A partir de ello, el BNE busca evaluar la dinámica de cada elemento en la cadena de valor del sector energético y poner a disposición del público información valiosa para el desarrollo del mismo.
Además, esta información es esencial para desarrollar los instrumentos de planeación del sector energético, evaluar su desempeño, así como para el diseño de políticas públicas y la toma de decisiones estratégicas alineados al Plan Nacional de Desarrollo (PND) 2025-2030 y al Programa Sectorial de Energía (PROSENER) 2025-2030.
El BNE se basa únicamente en un análisis energético, por lo que excluye aspectos económicos como el Producto Interno Bruto (PIB), los precios y las tarifas de la energía o la inflación. Esto se debe a que el documento ofrece un panorama general sobre el sector energético a partir de las bases energéticas objetivas con las que cuenta el país, con lo que se ofrece, no solo una visión del presente, sino una proyección del futuro energético de México.



\subsection{Fundamento jurídico del Balance Nacional de Energía}

El BNE 2024 se elabora en el marco de la Reforma Constitucional sobre áreas y empresas estratégicas y en materia de simplificación orgánica aprobadas en octubre y diciembre de 2024, la cual constituye un hito histórico para dar entrada a la nueva política energética del Estado mexicano. La reforma, que se articula en torno a la modificación de los artículos 25, 27 y 28 de la Constitución Política de los Estados Unidos Mexicanos (CPEUM)\footnote{Decreto por el que se reforman el párrafo quinto del artículo 25, los párrafos sexto y séptimo del artículo 27 y el párrafo cuarto del artículo 28 de la Constitución Política de los Estados Unidos Mexicanos, en materia de áreas y empresas estratégicas.}, se consolidó con la armonización de la legislación secundaria, que comprendió la expedición de ocho nuevas leyes \footnote{Ley de la Empresa Pública del Estado, Comisión Federal de Electricidad; Ley de la Empresa Pública del Estado, Petróleos Mexicanos; Ley del Sector Eléctrico; Ley del Sector Hidrocarburos; Ley de Planeación y Transición Energética; Ley de Biocombustibles; Ley de Geotermia y la Ley de la Comisión Nacional de Energía.} y la modificación a tres leyes asociadas\footnote{Ley de Ingresos sobre Hidrocarburos, Ley del Fondo Mexicano del Petróleo para la Estabilización y el Desarrollo y la Ley Orgánica de la Administración Pública Federal.} y una nota que, en su conjunto, fortalecen la rectoría del Estado en el sector energético.

Este nuevo orden normativo tiene como ejes fundamentales garantizar la seguridad y autosuficiencia energéticas, fortalecer la soberanía nacional, consolidar a la Comisión Federal de Electricidad (CFE) y Petróleos Mexicanos (PEMEX) como Empresas Públicas del Estado al servicio del pueblo, fortalecer a los demás organismos públicos del sector y robustecer las capacidades del Estado para ejercer una planeación vinculante, eficaz y estratégica. Asimismo, incorpora el principio de Justicia Energética que coloca en el centro de la acción gubernamental el acceso universal a energía asequible, segura y limpia para la atención de necesidades básicas, así como el combate frontal al comercio ilícito de combustibles y la transición energética sustentable, con una visión de largo plazo que armoniza el desarrollo económico con la protección del medio ambiente.

El papel rector y estratégico del Estado en el sector energético nacional, así como el aprovechamiento de los recursos naturales contenidos en el subsuelo (propiedad inalienable e imprescriptible de la Nación), se ejerce a través de las empresas públicas, sin que ello implique el establecimiento de monopolios, conforme lo disponen los artículos 27 y 28 constitucionales.

De acuerdo con el artículo 28 constitucional y con lo dispuesto en el artículo 2 de la Ley de Planeación y Transición Energética (LPTE), la SENER es la encargada de conducir la política energética nacional. Para ello, la SENER requiere elaborar documentos que coadyuven a su implementación, como el BNE, cuyo contenido se propone servir de marco de referencia para el cumplimiento de los objetivos propuestos en el PND 2025-2030 y el PROSENER 2025-2030.

Asimismo, conforme a lo estipulado en el artículo 35 de la LPTE, la SENER debe elaborar y publicar anualmente el BNE con escenarios prospectivos de energía que presente la información de origen y destino de la energía en el territorio nacional del año anterior.

Es importante mencionar que, derivado de los artículos 8, 20, 33 y 34 de la LPTE, a la SENER le corresponde implementar y administrar el Sistema Nacional de Información Energética (SNIEn), con el apoyo del Consejo de Planeación Energética, con el objeto de registrar, organizar, actualizar y difundir la información del sector energético. A su vez, le corresponde solicitar información a las Dependencias de la Administración Pública Federal, Empresas Públicas del Estado, gobiernos locales y otros organismos públicos o privados, así como a particulares, de acuerdo con lo establecido en el marco regulatorio del SNIEn.

Con respecto a la operación e implementación del SNIEn, la SENER debe emitir las disposiciones administrativas para la entrega y actualización de información, así como para la publicación de información del Sector Energético, de acuerdo con la política energética nacional. Además, las disposiciones permitirán la integración del SNIE a los diferentes sistemas de información del Sector Energético, y ayudará a la promoción de acuerdos con otras dependencias de la Administración Pública Federal, a fin de compartir información estratégica del Sector Energético, contratar licencias de software, infraestructura y servicios especializados, para la gestión protección y análisis de la información del SNIE.

Adicionalmente, el marco del SNIEn indica que la SENER debe observar las recomendaciones del Consejo y lo establecido en la Ley del Sistema Nacional de Información Estadística y Geografía, para la operación e implementación del Sistema. Además, la SENER debe publicar información e integrar los diferentes sistemas del Sector Energético de acuerdo a la política energética nacional.

Derivado de lo anterior, el BNE se encuentra aún en proceso de consolidación, en tanto el SNIEn continúa en desarrollo. Una vez que el Sistema se consolide, la calidad y cobertura de la información integrada en el BNE se verán fortalecidas. El SNIEn permitirá disponer de estadísticas energéticas más completas, consistentes y oportunas, lo que contribuirá también a una planeación más eficaz del sector energético.


\subsection{Objetivos del Balance Nacional de Energía}

El BNE 2024 es un documento cuyo objetivo general es llevar a cabo la evaluación sobre del sector energético de México, y proporcionar al público en general la información de toda la cadena de valor de la producción y consumo tanto de la energía primaria como secundaria. Sus objetivos específicos son los siguientes:

\begin{itemize}
  \item Servir de punto de partida para la planeación energética vinculante y la elaboración de los cuatro instrumentos de planeación del Sector Energético, los cuales deben estar alineados al PND 2025 - 2030: la Estrategia Nacional de Transición Energética, PLATEASE, PLADESE y PLADESHi.
  \item Evaluar y monitorear el funcionamiento del sector energético y el avance de la autosuficiencia, independencia, seguridad, transición, así como la planeación energética vinculante.
  \item Proporcionar información oficial energética pública, confiable, asequible, integral, verificable y didáctica, de acceso libre y gratuito para los interesados en el sector energético nacional.
  \item Establecerse como una fuente de información de referencia nacional e internacional para análisis especializados sobre el estado actual y futuro de la industria energética en México.
  \item Ofrecer información que posibilite el análisis del papel del sector en la transición energética baja en carbón y la lucha contra el cambio climático.
  \item Proporcionar proyecciones generales del sector en un horizonte de 15 años, como base para los cuatro instrumentos de planeación del Sector Energético.
  \item Garantizar que la información presentada esté alineada con las normativas internacionales y los estándares de calidad requeridos para la comparabilidad y transparencia.
\end{itemize}


\subsection{ Estructura del Balance Nacional de Energía}

El BNE 2024 se compone de seis capítulos. En el primero, se aborda el marco normativo y conceptual, así como los objetivos del Balance Nacional de Energía, mientras que en el capítulo dos se ahonda en la infraestructura energética, extracción, transporte, distribución, procesamiento y comercialización. Este capítulo tiene la finalidad de que el lector conozca el estado actual de la infraestructura nacional que permite la producción, así como el consumo de energía en sus distintas aplicaciones.

Además, en los capítulos tres y cuatro respectivamente, se muestra la producción y el resto de elementos que componen la oferta interna bruta, así como la transformación de la energía primaria, lo que da lugar a energéticos secundarios; el consumo final energético y no energético es revisado en el quinto capítulo. En el sexto, se presenta la información estadística para el Balance Nacional de Energía.

Adicionalmente, se incluye un glosario, las unidades y los factores de conversión para facilitar la interpretación de la información, así como una explicación sobre la metodología y las fuentes de información a partir de las cuales se reportan los datos publicados en el presente BNE.

Finalmente, se incluye un anexo estadístico con información histórica del sector energético que amplía el alcance del Balance, al ofrecer una visión retrospectiva de las variables analizadas en los capítulos dos al seis. Con ello, el BNE proporciona insumos que facilitan la elaboración de análisis especializados orientados a comprender el estado actual del Sector Energético Mexicano (SEM).


\portadaseccion{2}{Infraestructura energética}{}

\section{Infraestructura del Sector Hidrocarburos}




\subsection{Infraestructura del Sector Hidrocarburos}

México cuenta con infraestructura en el sector hidrocarburos que abarca toda la cadena de valor desde la exploración, producción, procesamiento, transporte, almacenamiento y distribución de los recursos petroleros y de gas natural. Construidas a lo largo de décadas, todas estas instalaciones actualmente operan bajo un esquema mixto donde participan tanto empresas públicas del Estado como actores privados, bajo un marco regulado por el Estado mexicano.
\begin{tabladoradoCorto}
  \caption{Tabla 2.10. Capacidad de almacenamiento y distribución de Petrolíferos, 2024}
  \label{tab:tabla_2_10__capacidad_de_almac}
  \begin{tabularx}{\textwidth}{Vvvv}
    \toprule
    \rowcolor{gobmxDorado} \encabezadodorado{MODALIDAD} & \encabezadodorado{TERMINALES} & \encabezadodorado{CAPACIDAD (MB)} & \encabezadodorado{} \\
    \midrule
    TAD de Pemex (Terrestres) & 62 & 9,097 & Capacidad operativa \\
    Terminales de operación portuaria & 10 & 4,768 & Capacidad operativa \\
    Terminales de operación marítima & 5 & 8,566 & Capacidad operativa \\
    Terminales de almacenamiento privadas & 15 & 12,772 & Capacidad operativa \\
    Almacenamiento en ASA & 52 & 685 & Capacidad nominal \\
    Almacenamiento de turbosina y gasavión de privados & 4  & 7 & Capacidad nominal \\
    Almacenamiento de turbosina GAFSACOMM & 9 & 149 & Capacidad nominal \\
    Terminales de almacenamiento de CFE & 49 & 10,276 & Capacidad disponible \\
    Distribuidores* & 247 & 1,204 & Capacidad nominal \\
    \bottomrule
  \end{tabularx}
\end{tabladoradoCorto}
\vspace{-4pt}
\fuente{Elaboración de SENER con información del Prontuario Estadístico de la SENER.}

\begin{tabladoradoCorto}
  \caption{Tabla 2.1 Reservas de Hidrocarburos de la Nación, 2024}
  \label{tab:tabla_2_1_reservas_de_hidrocar}
  \begin{tabularx}{\textwidth}{Vvvv}
    \toprule
    \rowcolor{gobmxDorado} \encabezadodorado{CATEGORIA} & \encabezadodorado{ACEITE 
(MMb)} & \encabezadodorado{GAS 
(MMMpc)} & \encabezadodorado{PCE 
(MMbpce)} \\
    \midrule
    Total 1p & 5,978 & 12,297 & 8,383 \\
    Total 2P & 11,078 & 23,302 & 15,530 \\
    Total 3P & 16,383 & 34,858 & 23,146 \\
    \bottomrule
  \end{tabularx}
\end{tabladoradoCorto}
\vspace{-4pt}
\fuente{Elaboración de SENER con información del Sistema de Información de Hidrocarburos de la SENER, Reservas al 01 de enero de 2024}



\subsubsection{Exploración y Extracción}

La infraestructura de exploración y extracción está destinada a la localización y obtención de recursos energéticos del subsuelo. Las actividades relacionadas con la exploración y extracción de hidrocarburos se concentran principalmente en las cuencas petroleras de la región marina somera del Golfo de México, así como en la porción terrestre de las cuencas del sureste y noreste del país.

El desarrollo en este sector depende del volumen de reservas de hidrocarburos disponibles, las cuales son certificadas anualmente según su certeza de recuperación. A pesar de que las reservas de hidrocarburos no se consideran parte de la infraestructura petrolera, su localización es el principal objetivo de la exploración petrolera y el punto de partida para la planeación y ejecución de proyectos de extracción.

Las cifras consolidadas al 1 de enero de 2024 para las reservas de aceite, gas, y petróleo crudo equivalente (PCE), agrupadas en las categorías 1P, 2P y 3P se presentan en la siguiente

Tabla 2.1

Es importante destacar que, pese a la disminución de las reservas de hidrocarburos en los últimos años, la tasa de restitución integral de las reservas probadas (1P) de PCE ha permanecido por encima del 100\% desde 2022 a 2024.

Las reservas permiten dimensionar el potencial productivo del país, además constituyen un factor estratégico fundamental que incentiva la inversión y promueve el desarrollo del sector energético. Con el objetivo de impulsar el crecimiento de la industria, el Estado Mexicano le otorgó a PEMEX y a empresas privadas los derechos para realizar actividades de exploración o extracción de hidrocarburos, a través de asignaciones y contratos. Las asignaciones fueron otorgadas de manera directa a PEMEX, mientras que los contratos se adjudicaron mediante procesos de licitación en las que participaron tanto empresas privadas como PEMEX. En este contexto, México cuenta con 412 asignaciones operadas por PEMEX, estas asignaciones se distribuyen de la siguiente manera: 314 son terrestres (8 comparten superficie tanto en tierra como en mar), 80 se ubican en aguas someras y 18 en aguas profundas. En cuanto a los contratos (principalmente privados), de los 104 vigentes, 51 se ubican en tierra y 53 son marinos, 30 de aguas someras y 23 de aguas profundas (Tabla 2.2).

Los pozos petroleros hacen posible la exploración, evaluación y producción de petróleo y gas. Durante enero de 2025 se contabilizaron 6,546 pozos operando, de los cuales 3,873 corresponden a pozos de petróleo y gas, y 2,673 a pozos de "gas no asociado. (Tabla 2.2)


\subsubsection{Transformación industrial}

La refinación del petróleo consiste en separar las distintas cadenas de hidrocarburos que lo componen mediante procesos como la destilación, o en modificarlas a través de técnicas como el craqueo o descomposición térmica y catalítica, con el fin de obtener productos específicos, tales como gasolinas, diesel, turbosina y otros derivados.

En México, el Sistema Nacional de Refinación está conformado por siete refinerías operadas por PEMEX, las cuales cuentan con una capacidad nominal total de procesamiento de 1,980 mil barriles diarios (Mbd) de petróleo crudo (Tabla 2.3.).

En línea con el objetivo de autosuficiencia energética, a inicios de 2022, con el propósito de fortalecer su participación en el mercado nacional de petrolíferos, PEMEX adquirió el 50\% restante y así obtener el total de la participación en Deer Park Refining Limited Partnership, la cual es una refinería ubicada en Texas, Estados Unidos. A pesar de que la refinería sea propiedad completamente de PEMEX, la producción que llega a México se debe considerar como importada y además tiene que cumplir con las reglas de importación, al ser productos manufacturados fuera del territorio nacional. No obstante, la refinería está produciendo para PEMEX y para los mercados donde tiene influencia, y representa valor a la empresa pública y al Estado mexicano.

En el caso del gas natural extraído del subsuelo, este contiene una mezcla de hidrocarburos que deben ser separados antes de su aprovechamiento. Este proceso se realiza en los Centros Procesadores de Gas (CPG), donde el gas es tratado para separar el metano de otros hidrocarburos más pesados. Como resultado, se obtiene gas seco y líquidos del gas natural, entre los que se encuentran el etano, el gas licuado de petróleo (GLP), las gasolinas naturales y las naftas.

En México operan nueve CPG administrados por PEMEX, los cuales cuentan con una capacidad instalada de 4,283 millones de pies cúbicos diarios (MMpcd) para el endulzamiento de gas, 135 MMpcd para el tratamiento de líquidos, 5,912 MMpcd para procesos criogénicos, y 560 mil barriles diarios (Mbd) para el fraccionamiento de líquidos (Tabla 2.4.).

La industria petroquímica desempeña un papel fundamental en la transformación de hidrocarburos como el gas natural y sus derivados. México contaba con ocho complejos petroquímicos; sin embargo, debido a políticas gubernamentales anteriores orientadas al abandono del sector y al deterioro de su infraestructura, dos de ellos fueron cerrados. Las razones argumentadas incluyeron la creciente falta de competitividad, el acceso limitado a mercados adecuados y la escasez de materia prima confiable. Estos factores afectaron la rentabilidad al grado de no cubrir ni sus costos variables, por lo que actualmente operan únicamente seis complejos petroquímicos, con una capacidad instalada total de 10,031 miles de toneladas (Mt) (véase Tabla 2.5.).

Frente a esta situación, la política energética vigente ha reorientado sus esfuerzos hacia la recuperación del sector petroquímico como componente central para fortalecer la soberanía energética nacional. En ese sentido, se han diseñado programas de mantenimiento integral y modernización operativa con el objetivo de revitalizar los complejos existentes y consolidar su viabilidad económica.


\subsubsection{Transporte, almacenamiento y distribución}

El gas seco obtenido de los CPG, así como el extraído directo de campos y el importado, puede ser transportado y distribuido hacia los usuarios finales a través de ductos. En 2024 la infraestructura de transporte de gas natural por dicho medio, cuenta con una longitud de 21,149 km, de los cuales 9,732 km (46\%) se encuentran interconectados e integrados al Sistema de Transporte y Almacenamiento Nacional Integrado de Gas Natural (SISTRANGAS), el cual es gestionado por el Centro Nacional de Control del Gas Natural (CENAGAS). Asimismo, el CENAGAS opera y da mantenimiento al Sistema Naco-Hermosillo (SNH), que consta de 355 km. El resto de los gasoductos, son operados y gestionados de forma independiente por privados.

En el SISTRANGAS, de los 9,732 km que lo conforman, 8,386 km (86.2\%) forman parte del Sistema Nacional de Gasoductos (SNG) y son propiedad del Estado, en tanto que 1,346 km de gasoductos son privados (13.8\%) y están interconectados e integrados al SISTRANGAS (Tabla 2.7.).

Es importante resaltar que el gas seco que se consume en el país se obtiene principalmente por importación. Al cierre de 2024, para la importación de este energético por ducto, se contaban con 24 puntos de internación con la frontera de Estados Unidos. En cuanto al almacenamiento de gas natural, existen 5 terminales de regasificación y almacenamiento de gas natural licuado (ver Tabla 2.8.).

Mientras que los petrolíferos producidos en las refinerías y los importados a través de las terminales de operación marítima y portuaria, son transportados por ductos, auto tanques, carro tanques o buque tanques para su almacenamiento. En 2024, el transporte de petrolíferos, que considera poliductos, ramales, ductos bidireccionales, entre otros, sumaron una extensión total de 9,582 km y una capacidad de 4,586 Mb. En el caso de GLP, la red existente de 4 ductos en el país suman 2,044 km (Tabla 2.9.).

Las Terminales de Almacenamiento y Distribución (TAD) reciben los petrolíferos desde los lugares de producción u origen, para almacenarlos y posteriormente su reparto. El transporte de los petrolíferos desde las TAD para su distribución a los puntos de expendio al público se realiza mediante autotanques propiedad de PEMEX o de empresas privadas. Esto permite suministrar petrolíferos a las estaciones de servicio distribuidas a lo largo y ancho del país. En la Tabla 2.10 se muestran las capacidades de almacenamiento y distribución de petrolíferos.

Con la finalidad de contar con instalaciones o sistemas, para conservar en depósito y resguardar el GLP propiedad de terceros, en 2024 se contaba con Plantas de Almacenamiento en México con una capacidad de 6,021 Mb. Además, para adquirir, recibir, guardar y conducir GLP a granel, se cuenta con plantas de distribución para repartición, traslado y enajenación a uno o varios destinos, con una capacidad en 2024 de 2,439 Mb (Ver Tabla 2.11.).


\subsection{Infraestructura del Sector Eléctrico}

La infraestructura del Sistema Eléctrico Nacional (SEN) cumple con dos funciones principales que permiten satisfacer la demanda de electricidad en el país: la generación de energía eléctrica y su transmisión y distribución a los centros de consumo. En el país se cuenta con más de 100 GW de capacidad de generación permisionada y más de 100 mil kilómetros en la Red Nacional de Transmisión (RNT), que conectan las principales fuentes de generación con la infraestructura de transformación (193 mil MVA). Asimismo, existen cerca de 900 mil kilómetros en las Redes Generales de Distribución (RGD) que se encargan de llevar la electricidad directamente a hogares, comercios e industrias.

En las siguientes secciones, se describen estos componentes para entender su importancia e integración en la infraestructura del sector eléctrico en el país.
\begin{tabladoradoCorto}
  \caption{Tabla 2.2 Asignaciones, contratos y pozos operando por ubicación, 2024}
  \label{tab:tabla_2_2_asignaciones__contra}
  \begin{tabularx}{\textwidth}{Vvvv}
    \toprule
    \rowcolor{gobmxDorado} \encabezadodorado{UBICACIÓN} & \encabezadodorado{ASIGNACIONES} & \encabezadodorado{CONTRATOS} & \encabezadodorado{POZOS OPERANDO} \\
    \midrule
    Terrestre & 314 & 51 & 5,924 \\
    Marinos* & 98 & 53 & 622 \\
    Total & 412 & 104 & 6,546 \\
    \bottomrule
  \end{tabularx}
\end{tabladoradoCorto}
\vspace{-4pt}
\fuente{Elaboración de SENER con información del Sistema de Información de Hidrocarburos, Rondas México y Portal de Asignaciones de la SENER}



\subsubsection{Capacidad de Generación de Energía Electrica}

La generación de electricidad sucede primordialmente a gran escala en centrales que utilizan diversas tecnologías. A lo largo del territorio nacional, existen 101,176 MW de capacidad permisionada en operación a 2024, distribuidos como se muestra en la Figura 2.2.

Como se observa en la Tabla 2.12., la tecnología de Ciclo Combinado es la que actualmente tiene una mayor aportación a la capacidad total del SEN, con cerca de 41 GW; seguida de la Térmica Convencional con 15 GW y la Hidroeléctrica con 13 GW aproximadamente.


\subsubsection{Transmisión y Disitribución}

Comparada con una extensa red de carreteras que transporta la electricidad desde las plantas de generación eléctrica, la RNT está integrada por más de 100 mil kilómetros de líneas de alta tensión que recorren el territorio del país. La Tabla 2.13. describe los dos grandes grupos de niveles de tensión para el año 2024: el primero, dedicado a la transmisión de entre 161 a 400 kV con 56,409 km, donde destacan las líneas de 400 kV (26,135 km) y 230 kV (29,743 km) para el transporte de energía eléctrica a gran escala; el segundo grupo, de 69 a 138 kV con 54,437 km, donde predominan líneas de 115 kV (48,725 km) que son vías primarias para la distribución regional y la interconexión de subestaciones. Esta infraestructura de 110,846 km, refleja la columna vertebral del sistema eléctrico mexicano.

Entre la RNT y la distribución a centros de consumo, se encuentra la infraestructura de transformación, es decir, los equipos de acondicionamiento de la corriente transportada según las necesidades de la infraestructura de transporte. Para 2024, se cuenta con una capacidad de 113,688 MVA en bancos de transformación. La Tabla 2.14 muestra la capacidad de transformación en la RNT y las RGD del MEM (subestaciones de alta a media tensión que forman parte del Mercado Eléctrico Mayorista).

Por otra parte, las RGD que entregan la electricidad a los usuarios finales, se componen de líneas de media tensión (2.4 - 34 kV), baja tensión (menos de 2.4 kV), así como de capacidad de transformación. En la Tabla 2.15 se muestra la infraestructura de distribución de energía eléctrica registrada en 2024.


\subsubsection{Capacidad Instalada de Generación Distribuida}

De conformidad con el artículo 3, fracción XXII de la Ley del Sector Eléctrico (LSE), la Generación Distribuida (GD) es la generación que se realiza por una Generadora Exenta y que se encuentra interconectada a un circuito de distribución que contenga una alta concentración de Centros de Carga. Cabe señalar, que el artículo 19 de la LSE establece que todas las Centrales Eléctricas con capacidad menor a 0.7 MW son Generadoras Exentas. En otras palabras, la GD es la generación de energía eléctrica a través de centrales eléctricas de pequeña escala (Generadoras Exentas), que no requieren ni cuentan con un permiso de generación otorgado por la autoridad reguladora, pero puede estar interconectada a un circuito de distribución. Generalmente, estas centrales se encuentran cerca del punto de consumo y abastecen parcial o totalmente la demanda local.

Al cierre de 2024, la Comisión Reguladora de Energía registró un crecimiento de la capacidad instalada de la GD cercana al 30\% con respecto al año anterior, resultando en 4,447.92 MW, bajo las figuras de Generador Exento de la anterior Ley de la Industria Eléctrica (LIE), y los Contratos de Interconexión de Pequeña y Mediana Escala, de la la antigua Ley del Servicio Público de Energía Eléctrica (LSPEE) (Tabla 2.16).

Si bien, cualquier tipo de tecnología de generación puede ser aplicada a una Central Eléctrica de GD, la capacidad que actualmente se encuentra instalada predominantemente es de tecnología solar fotovoltaica, con una participación de 99.41\% del total. Lo que deja al resto de las tecnologías (eólica, bioenergía y otras) con una participación de apenas el 0.59\% en su conjunto. Como se puede observar en la Figura 2.3, Jalisco y Nuevo León son las entidades federativas que concentran cerca del 25\% de la capacidad instalada en la República Mexicana.


\subsection{Infraestructura de consumo final y aprovechamiento directo}

La infraestructura de consumo final está compuesta por el conjunto de equipos, vehículos, edificaciones, instalaciones y sistemas que utilizan energéticos para realizar funciones productivas, de transporte, comerciales, residenciales o de servicios. Esta composición es reflejo de la diversidad sectorial y territorial del país, por lo que su caracterización es clave para el análisis de los patrones de consumo.

En el sector transporte, que representa el mayor componente del consumo final energético nacional (45.59\% en 2024), la modalidad define la participación que tendrá en el consumo de energía. El transporte terrestre (automóviles particulares, taxis, motocicletas, camionetas de reparto, autobuses de pasajeros y camiones de carga), está dominado por un parque vehicular en circulación que se estima cercano a 50 millones de unidades. La infraestructura ferroviaria, aunque más limitada, cuenta con más de 1,300 locomotoras en operación, con más de 160 vagones para el transporte de pasajeros y más de 30 mil carros para carga (incluyendo plataformas, tolvas y vagones especializados). El transporte marítimo y fluvial utiliza embarcaciones de carga, ferrys y lanchas de transporte local en ríos y costas. En el transporte aéreo, el consumo energético se concentra en aeronaves comerciales y de carga que operan con turbosina, con una flota nacional de aproximadamente 5,200 aeronaves activas.

El sector industrial contaba con más de 640 mil unidades económicas, en 2024, con patrones de consumo altamente diferenciados según la rama productiva. La minería, cementera, siderurgia y metalurgia y la producción de papel se caracterizan por utilizar maquinaria pesada, hornos, calderas, intercambiadores de calor, sistemas de bombeo, ventilación y compresión, así como infraestructura para climatización e iluminación. En cambio, otras industrias como la química, alimentaria, textil o automotriz concentran su demanda en operaciones de maquinado, refrigeración, climatización y procesos específicos.

El sector comercio y servicios integra una infraestructura energética amplia, debido a la cantidad de unidades económicas existentes (más de 2,600,000 en 2024). Además, grandes edificios corporativos, instituciones de salud y educación, y oficinas de la administración pública, requieren un suministro energético constante para mantener condiciones adecuadas de operación, seguridad y confort.

El consumo en el sector residencial está directamente ligado al número de viviendas, la cantidad de habitantes por hogar y la penetración de electrodomésticos, tecnologías y dispositivos para satisfacer necesidades energéticas y confort. Según datos del INEGI, en 2020, de las más de 33 millones de viviendas particulares habitadas, casi el 99\% contó con suministro de energía eléctrica; y un promedio de 66\% contó con diversos electrodomésticos. En los hogares, la infraestructura de consumo incluye una vasta gama de dispositivos eléctricos (refrigeradores, lavadoras, televisores, equipos de cómputo), y equipamiento para iluminación, calentamiento de agua, cocción de alimentos y climatización (aires acondicionados y calefactores de espacios).

En cuanto al sector público, existen dos tipos de infraestructura energética relevantes. El alumbrado público, con más de 9 millones de luminarias en funcionamiento a lo largo del territorio nacional, representa un componente esencial del consumo eléctrico municipal. En paralelo, los sistemas de bombeo y tratamiento de agua potable y residual dependen de la electricidad para su funcionamiento continuo. Estos sistemas son operados principalmente por los más de 2,500 organismos públicos; aunque también existen cerca de 2,270 prestadores de servicio de agua de gestión social o privada.

Por su parte, el sector agropecuario utiliza una variada infraestructura de consumo esencial para la producción de alimentos. Entre los equipos destacan los tractores, cosechadoras y sembradoras y sistemas de riego. Mientras que para los procesos de post-cosecha, se usan secadores, molinos y equipos de refrigeración para almacenamiento de alimentos. Para la industria producción pecuaria, se utilizan equipos de ventilación, sistemas de alimentación automática, ordeñadoras eléctricas y calentadores.

Además, en México, los sistemas de aprovechamiento energético directo representan una alternativa eficiente y adecuada para diversos sectores productivos. Entre ellos, destacan particularmente los calentadores solares de agua (CSA), cuya adopción ha crecido significativamente en los últimos años. Según datos de Fabricantes Mexicanos de Energías Renovables, A.C. a 2024, el país cuenta con una capacidad instalada mayor a los 4,800 MW térmicos que representa un crecimiento del 34\% con respecto a lo instalado en el 2020 (SECIHTI, 2022). Esta tecnología se implementa principalmente en viviendas, aunque también ha tenido un crecimiento importante en establecimientos hoteleros, centros hospitalarios, gimnasios y procesos industriales, durante los últimos 15 años.
\begin{tabladoradoCorto}
  \caption{Tabla 2.3 Capacidad de refinación de petróleo crudo en México, 2024}
  \label{tab:tabla_2_3_capacidad_de_refinac}
  \begin{tabularx}{\textwidth}{Vvv}
    \toprule
    \rowcolor{gobmxDorado} \encabezadodorado{REFINERÍA} & \encabezadodorado{UBICACIÓN} & \encabezadodorado{Mbd} \\
    \midrule
    Ing. Héctor R. Lara Sosa & Cadereyta, Nuevo León & 275 \\
    Francisco l. Madero & Ciudad Madero, Tamaulipas & 190 \\
    Refinería Olmeca & Dos Bocas, Tabasco & 340 \\
    General Lázaro Cárdenas & Minatitlán, Veracruz & 285 \\
    Ing. Antonio M. Amor & Salamanca, Guanajuato & 245 \\
    Ing. Antonio Dovalí Jaime & Salina Cruz, Oaxaca & 330 \\
    Refinería Miguel Hidalgo & Tula, Hidalgo & 315 \\
    Sistema Nacional de Refinación &  & 1,980 \\
    \bottomrule
  \end{tabularx}
\end{tabladoradoCorto}
\vspace{-4pt}
\fuente{Elaboración de SENER con información del Anuario Estadístico de PEMEX 2024.}



\portadaseccion{3}{}{}

\section{Oferta interna bruta}




\subsection{Producción de energía primaria}




\subsection{Comercio exterior de energía}




\subsection{Variación de inventarios}




\subsection{Oferta total}




\subsection{Energía no aprovechada}




\subsection{Oferta interna bruta}




\portadaseccion{4}{}{}

\section{Transformaciones, Consumo Propio y Pérdidas}




\subsection{Refinación}




\subsection{Procesamiento de gas}




\subsection{Generación eléctrica}




\subsection{Consumo propio del sector}




\subsection{Pérdidas}




\portadaseccion{5}{}{}

\section{Consumo final}




\subsection{Sector agropecuario}




\subsection{Sector industrial}




\subsection{Sector comercial y servicios}




\subsection{Sector residencial}




\subsection{Sector público}




\subsection{Sector transporte}




\subsection{Consumo no energético}




\portadaseccion{6}{}{}

\section{Balance Nacional de Energía }




\subsection{Principales cuentas del Balance Nacional de Energía}


\begin{tabladoradoLargo}
  \begin{xltabular}{\textwidth}{YZZZZZ}
    \caption{Tabla 6.1. Matriz origen destino oferta interna bruta y transformaciones de energéticos primarios, 2024}\label{tab:tabla_6_1__matriz_origen_desti}\\
    \toprule
    \rowcolor{gobmxDorado} \encabezadodorado{} & \encabezadodorado{Carbón mineral} & \encabezadodorado{Petróleo crudo} & \encabezadodorado{Condensados} & \encabezadodorado{Gas natural} & \encabezadodorado{Energía Nuclear} \\
    \midrule
    \endfirsthead

    \multicolumn{6}{l}{\small\textit{Continuación...}} \\
    \toprule
    \rowcolor{gobmxDorado} \encabezadodorado{} & \encabezadodorado{Carbón mineral} & \encabezadodorado{Petróleo crudo} & \encabezadodorado{Condensados} & \encabezadodorado{Gas natural} & \encabezadodorado{Energía Nuclear} \\
    \midrule
    \endhead

    \midrule
    \multicolumn{6}{r}{\small\textit{Continúa en la siguiente página...}} \\
    \endfoot

    \bottomrule
    \endlastfoot

    Oferta interna bruta & 219.46 & 1,747.67 & 636.74 & 1,648.09 & 129.51 \\
    \hspace{1.35em} Oferta total & 219.56 & 3,674.19 & 637.14 & 1,853.93 & 129.51 \\
    \hspace{2.7em} Producción & 14.99 & 3,654 & 638 & 1,868 & 129.51 \\
    \hspace{2.7em} Importación & 206.15 & 13.2 & 0 & 0 & 0 \\
    \hspace{2.7em} Variación de inventarios & -1.58 & 6.98 & -0.86 & -14.07 & 0 \\
    \hspace{1.35em} Exportación & -0.09 & -1,873 & 0 & 0 & 0 \\
    \hspace{1.35em} Energía no aprovechada & 0 & -53.52 & -0.4 & -205.84 & 0 \\
    Total transformaciónes & -147.47 & -1,793.27 & -638.37 & -1,052.4 & -129.51 \\
    \hspace{1.35em} Coquizadoras y hornos & 0 & 0 & 0 & 0 & 0 \\
    \hspace{1.35em} Refinerías y despuntadoras & 0 & -1,793.27 & -326.21 & 0 & 0 \\
    \hspace{1.35em} Plantas de gas y fraccionadoras & 0 & 0 & -312.16 & -1,052.4 & 0 \\
    \hspace{1.35em} Centrales eléctricas & -147.47 & 0 & 0 & 0 & -129.51 \\
    \hspace{2.7em} Carboeléctrica & -147.47 & 0 & 0 & 0 & 0 \\
    \hspace{2.7em} Termica convencional & 0 & 0 & 0 & 0 & 0 \\
    \hspace{2.7em} Combustión interna & 0 & 0 & 0 & 0 & 0 \\
    \hspace{2.7em} Turbogas & 0 & 0 & 0 & 0 & 0 \\
    \hspace{2.7em} Ciclo combinado & 0 & 0 & 0 & 0 & 0 \\
    \hspace{2.7em} Nucleoeléctrica & 0 & 0 & 0 & 0 & -129.51 \\
    \hspace{2.7em} Cogeneración & 0 & 0 & 0 & 0 & 0 \\
    \hspace{2.7em} Hidroeléctrica & 0 & 0 & 0 & 0 & 0 \\
    \hspace{2.7em} Geotérmica & 0 & 0 & 0 & 0 & 0 \\
    \hspace{2.7em} Eólica & 0 & 0 & 0 & 0 & 0 \\
    \hspace{2.7em} Solar fotovoltáica & 0 & 0 & 0 & 0 & 0 \\
  \end{xltabular}

  \vspace{0.25em}
  {\small\textit{Continuación Tabla. Tabla 6.1. Matriz origen destino oferta interna bruta y transformaciones de energéticos primarios, 2024}}
  \vspace{0.15em}

  \begin{xltabular}{\textwidth}{YZZZZZ}
    \toprule
    \rowcolor{gobmxDorado} \encabezadodorado{} & \encabezadodorado{Energia Hidraúlica} & \encabezadodorado{Geoenergía} & \encabezadodorado{Energía solar} & \encabezadodorado{Energía eólica} & \encabezadodorado{Bagazo de caña} \\
    \midrule
    \endfirsthead

    \multicolumn{6}{l}{\small\textit{Continuación...}} \\
    \toprule
    \rowcolor{gobmxDorado} \encabezadodorado{} & \encabezadodorado{Energia Hidraúlica} & \encabezadodorado{Geoenergía} & \encabezadodorado{Energía solar} & \encabezadodorado{Energía eólica} & \encabezadodorado{Bagazo de caña} \\
    \midrule
    \endhead

    \midrule
    \multicolumn{6}{r}{\small\textit{Continúa en la siguiente página...}} \\
    \endfoot

    \bottomrule
    \endlastfoot

    Oferta interna bruta & 262.18 & 91.32 & 252.18 & 220.18 & 94.72 \\
    \hspace{1.35em} Oferta total & 262.18 & 91.32 & 252.18 & 220.18 & 95.77 \\
    \hspace{2.7em} Producción & 262.18 & 91.32 & 252.18 & 220.18 & 95.77 \\
    \hspace{2.7em} Importación & 0 & 0 & 0 & 0 & 0 \\
    \hspace{2.7em} Variación de inventarios & 0 & 0 & 0 & 0 & 0 \\
    \hspace{1.35em} Exportación & 0 & 0 & 0 & 0 & 0 \\
    \hspace{1.35em} Energía no aprovechada & 0 & 0 & 0 & 0 & -1.05 \\
    Total transformaciónes & -262.18 & -91.32 & -205.7 & -220.18 & -61.5 \\
    \hspace{1.35em} Coquizadoras y hornos & 0 & 0 & 0 & 0 & 0 \\
    \hspace{1.35em} Refinerías y despuntadoras & 0 & 0 & 0 & 0 & 0 \\
    \hspace{1.35em} Plantas de gas y fraccionadoras & 0 & 0 & 0 & 0 & 0 \\
    \hspace{1.35em} Centrales eléctricas & -262.18 & -91.32 & -205.7 & -220.18 & -61.5 \\
    \hspace{2.7em} Carboeléctrica & 0 & 0 & 0 & 0 & 0 \\
    \hspace{2.7em} Termica convencional & 0 & 0 & 0 & 0 & -56.67 \\
    \hspace{2.7em} Combustión interna & 0 & 0 & 0 & 0 & -4.83 \\
    \hspace{2.7em} Turbogas & 0 & 0 & 0 & 0 & 0 \\
    \hspace{2.7em} Ciclo combinado & 0 & 0 & 0 & 0 & 0 \\
    \hspace{2.7em} Nucleoeléctrica & 0 & 0 & 0 & 0 & 0 \\
    \hspace{2.7em} Cogeneración & 0 & 0 & 0 & 0 & 0 \\
    \hspace{2.7em} Hidroeléctrica & -262.18 & 0 & 0 & 0 & 0 \\
    \hspace{2.7em} Geotérmica & 0 & -91.32 & 0 & 0 & 0 \\
    \hspace{2.7em} Eólica & 0 & 0 & 0 & -220.18 & 0 \\
    \hspace{2.7em} Solar fotovoltáica & 0 & 0 & -205.7 & 0 & 0 \\
  \end{xltabular}

  \vspace{0.25em}
  {\small\textit{Continuación Tabla. Tabla 6.1. Matriz origen destino oferta interna bruta y transformaciones de energéticos primarios, 2024}}
  \vspace{0.15em}

  \begin{xltabular}{\textwidth}{YZZZ}
    \toprule
    \rowcolor{gobmxDorado} \encabezadodorado{} & \encabezadodorado{Leña} & \encabezadodorado{Biogás} & \encabezadodorado{Total energéticos primarios} \\
    \midrule
    \endfirsthead

    \multicolumn{4}{l}{\small\textit{Continuación...}} \\
    \toprule
    \rowcolor{gobmxDorado} \encabezadodorado{} & \encabezadodorado{Leña} & \encabezadodorado{Biogás} & \encabezadodorado{Total energéticos primarios} \\
    \midrule
    \endhead

    \midrule
    \multicolumn{4}{r}{\small\textit{Continúa en la siguiente página...}} \\
    \endfoot

    \bottomrule
    \endlastfoot

    Oferta interna bruta & 129.42 & 3.25 & 5,434.71 \\
    \hspace{1.35em} Oferta total & 129.42 & 3.25 & 7,566.61 \\
    \hspace{2.7em} Producción & 129.42 & 3.25 & 7,358.78 \\
    \hspace{2.7em} Importación & 0 & 0 & 219.35 \\
    \hspace{2.7em} Variación de inventarios & 0 & 0 & -9.52 \\
    \hspace{1.35em} Exportación & 0 & 0 & -1,873.09 \\
    \hspace{1.35em} Energía no aprovechada & 0 & 0 & -260.8 \\
    Total transformaciónes & 0 & -3.25 & -4,605.13 \\
    \hspace{1.35em} Coquizadoras y hornos & 0 & 0 & 0 \\
    \hspace{1.35em} Refinerías y despuntadoras & 0 & 0 & -2,119.48 \\
    \hspace{1.35em} Plantas de gas y fraccionadoras & 0 & 0 & -1,364.56 \\
    \hspace{1.35em} Centrales eléctricas & 0 & -3.25 & -1,121.09 \\
    \hspace{2.7em} Carboeléctrica & 0 & 0 & -147.47 \\
    \hspace{2.7em} Termica convencional & 0 & -0.42 & -57.08 \\
    \hspace{2.7em} Combustión interna & 0 & 0 & -4.83 \\
    \hspace{2.7em} Turbogas & 0 & -2.83 & -2.83 \\
    \hspace{2.7em} Ciclo combinado & 0 & 0 & 0 \\
    \hspace{2.7em} Nucleoeléctrica & 0 & 0 & -129.51 \\
    \hspace{2.7em} Cogeneración & 0 & 0 & 0 \\
    \hspace{2.7em} Hidroeléctrica & 0 & 0 & -262.18 \\
    \hspace{2.7em} Geotérmica & 0 & 0 & -91.32 \\
    \hspace{2.7em} Eólica & 0 & 0 & -220.18 \\
    \hspace{2.7em} Solar fotovoltáica & 0 & 0 & -205.7 \\
  \end{xltabular}
\end{tabladoradoLargo}
\vspace{-4pt}
\fuente{Elaboración SENER}



\subsection{Matriz del Balance Nacional de Energía}


\begin{tabladoradoLargo}
  \begin{xltabular}{\textwidth}{YZZZZZ}
    \caption{Tabla 6.2. Matriz origen destino, consumo propio, pérdidas y consumo final de energéticos primarios, 2024}\label{tab:tabla_6_2__matriz_origen_desti}\\
    \toprule
    \rowcolor{gobmxDorado} \encabezadodorado{} & \encabezadodorado{Carbón mineral} & \encabezadodorado{Petróleo crudo} & \encabezadodorado{Condensados} & \encabezadodorado{Gas natural} & \encabezadodorado{Energía Nuclear} \\
    \midrule
    \endfirsthead

    \multicolumn{6}{l}{\small\textit{Continuación...}} \\
    \toprule
    \rowcolor{gobmxDorado} \encabezadodorado{} & \encabezadodorado{Carbón mineral} & \encabezadodorado{Petróleo crudo} & \encabezadodorado{Condensados} & \encabezadodorado{Gas natural} & \encabezadodorado{Energía Nuclear} \\
    \midrule
    \endhead

    \midrule
    \multicolumn{6}{r}{\small\textit{Continúa en la siguiente página...}} \\
    \endfoot

    \bottomrule
    \endlastfoot

    Consumo del sector y pérdidas & 0 & 45.6 & 1.63 & -595.69 & 0 \\
    \hspace{1.35em} Consumo propio del sector & 0 & -12.4 & 0 & -595.69 & 0 \\
    \hspace{1.35em} Diferencia estadística & 0 & 69.82 & 1.63 & 0 & 0 \\
    \hspace{1.35em} Pérdidas técnicas por transporte, transmisión y distribución & 0 & -11.82 & 0 & 0 & 0 \\
    \hspace{2.7em} Pérdidas en transporte y transmisión & 0 & -11.82 & 0 & 0 & 0 \\
    \hspace{2.7em} Pérdidas técnicas en distribución & 0 & 0 & 0 & 0 & 0 \\
    \hspace{2.7em} Pérdidas no técnicas  en distribución & 0 & 0 & 0 & 0 & 0 \\
    Consumo final total & 72 & 0 & 0 & 0 & 0 \\
    \hspace{1.35em} Consumo final no energético & 0 & 0 & 0 & 0 & 0 \\
    \hspace{2.7em} Petroquímica Pemex & 0 & 0 & 0 & 0 & 0 \\
    \hspace{2.7em} Otras ramas económicas & 0 & 0 & 0 & 0 & 0 \\
    Consumo final energético & 72 & 0 & 0 & 0 & 0 \\
    \hspace{2.7em} Agropecuario & 0 & 0 & 0 & 0 & 0 \\
    \hspace{2.7em} Industrial & 72 & 0 & 0 & 0 & 0 \\
    \hspace{2.7em} Comercial y servicios & 0 & 0 & 0 & 0 & 0 \\
    \hspace{2.7em} Residencial & 0 & 0 & 0 & 0 & 0 \\
    \hspace{2.7em} Público & 0 & 0 & 0 & 0 & 0 \\
    \hspace{2.7em} Transporte & 0 & 0 & 0 & 0 & 0 \\
  \end{xltabular}

  \vspace{0.25em}
  {\small\textit{Continuación Tabla. Tabla 6.2. Matriz origen destino, consumo propio, pérdidas y consumo final de energéticos primarios, 2024}}
  \vspace{0.15em}

  \begin{xltabular}{\textwidth}{YZZZZZ}
    \toprule
    \rowcolor{gobmxDorado} \encabezadodorado{} & \encabezadodorado{Energia Hidraúlica} & \encabezadodorado{Geoenergía} & \encabezadodorado{Energía solar} & \encabezadodorado{Energía eólica} & \encabezadodorado{Bagazo de caña} \\
    \midrule
    \endfirsthead

    \multicolumn{6}{l}{\small\textit{Continuación...}} \\
    \toprule
    \rowcolor{gobmxDorado} \encabezadodorado{} & \encabezadodorado{Energia Hidraúlica} & \encabezadodorado{Geoenergía} & \encabezadodorado{Energía solar} & \encabezadodorado{Energía eólica} & \encabezadodorado{Bagazo de caña} \\
    \midrule
    \endhead

    \midrule
    \multicolumn{6}{r}{\small\textit{Continúa en la siguiente página...}} \\
    \endfoot

    \bottomrule
    \endlastfoot

    Consumo del sector y pérdidas & 0 & 0 & 0 & 0 & 0 \\
    \hspace{1.35em} Consumo propio del sector & 0 & 0 & 0 & 0 & 0 \\
    \hspace{1.35em} Diferencia estadística & 0 & 0 & 0 & 0 & 0 \\
    \hspace{1.35em} Pérdidas técnicas por transporte, transmisión y distribución & 0 & 0 & 0 & 0 & 0 \\
    \hspace{2.7em} Pérdidas en transporte y transmisión & 0 & 0 & 0 & 0 & 0 \\
    \hspace{2.7em} Pérdidas técnicas en distribución & 0 & 0 & 0 & 0 & 0 \\
    \hspace{2.7em} Pérdidas no técnicas  en distribución & 0 & 0 & 0 & 0 & 0 \\
    Consumo final total & 0 & 0 & 46.48 & 0 & 33.22 \\
    \hspace{1.35em} Consumo final no energético & 0 & 0 & 0 & 0 & 0.04 \\
    \hspace{2.7em} Petroquímica Pemex & 0 & 0 & 0 & 0 & 0 \\
    \hspace{2.7em} Otras ramas económicas & 0 & 0 & 0 & 0 & 0.04 \\
    Consumo final energético & 0 & 0 & 46.48 & 0 & 33.18 \\
    \hspace{2.7em} Agropecuario & 0 & 0 & 0 & 0 & 0 \\
    \hspace{2.7em} Industrial & 0 & 0 & 31.49 & 0 & 33.18 \\
    \hspace{2.7em} Comercial y servicios & 0 & 0 & 4.41 & 0 & 0 \\
    \hspace{2.7em} Residencial & 0 & 0 & 10.58 & 0 & 0 \\
    \hspace{2.7em} Público & 0 & 0 & 0 & 0 & 0 \\
    \hspace{2.7em} Transporte & 0 & 0 & 0 & 0 & 0 \\
  \end{xltabular}

  \vspace{0.25em}
  {\small\textit{Continuación Tabla. Tabla 6.2. Matriz origen destino, consumo propio, pérdidas y consumo final de energéticos primarios, 2024}}
  \vspace{0.15em}

  \begin{xltabular}{\textwidth}{YZZZ}
    \toprule
    \rowcolor{gobmxDorado} \encabezadodorado{} & \encabezadodorado{Leña} & \encabezadodorado{Biogás} & \encabezadodorado{Total energéticos primarios} \\
    \midrule
    \endfirsthead

    \multicolumn{4}{l}{\small\textit{Continuación...}} \\
    \toprule
    \rowcolor{gobmxDorado} \encabezadodorado{} & \encabezadodorado{Leña} & \encabezadodorado{Biogás} & \encabezadodorado{Total energéticos primarios} \\
    \midrule
    \endhead

    \midrule
    \multicolumn{4}{r}{\small\textit{Continúa en la siguiente página...}} \\
    \endfoot

    \bottomrule
    \endlastfoot

    Consumo del sector y pérdidas & 0 & 0 & -548.46 \\
    \hspace{1.35em} Consumo propio del sector & 0 & 0 & -608.09 \\
    \hspace{1.35em} Diferencia estadística & 0 & 0 & 71.45 \\
    \hspace{1.35em} Pérdidas técnicas por transporte, transmisión y distribución & 0 & 0 & -11.82 \\
    \hspace{2.7em} Pérdidas en transporte y transmisión & 0 & 0 & -11.82 \\
    \hspace{2.7em} Pérdidas técnicas en distribución & 0 & 0 & 0 \\
    \hspace{2.7em} Pérdidas no técnicas  en distribución & 0 & 0 & 0 \\
    Consumo final total & 129.42 & 0 & 281.12 \\
    \hspace{1.35em} Consumo final no energético & 0 & 0 & 0.04 \\
    \hspace{2.7em} Petroquímica Pemex & 0 & 0 & 0 \\
    \hspace{2.7em} Otras ramas económicas & 0 & 0 & 0.04 \\
    Consumo final energético & 129.42 & 0 & 281.08 \\
    \hspace{2.7em} Agropecuario & 0 & 0 & 0 \\
    \hspace{2.7em} Industrial & 9.82 & 0 & 146.49 \\
    \hspace{2.7em} Comercial y servicios & 4.66 & 0 & 9.07 \\
    \hspace{2.7em} Residencial & 114.93 & 0 & 125.51 \\
    \hspace{2.7em} Público & 0 & 0 & 0 \\
    \hspace{2.7em} Transporte & 0 & 0 & 0 \\
  \end{xltabular}
\end{tabladoradoLargo}
\vspace{-4pt}
\fuente{Elaboración SENER}



\subsection{Principales flujos del Balance Nacional de Energía}


\begin{tabladoradoLargo}
  \begin{xltabular}{\textwidth}{YZZZZZ}
    \caption{Tabla 6.3. Matriz origen destino, oferta interna bruta y transformaciones de energéticos secundarios, 2024}\label{tab:tabla_6_3__matriz_origen_desti}\\
    \toprule
    \rowcolor{gobmxDorado} \encabezadodorado{} & \encabezadodorado{Coque de carbón} & \encabezadodorado{Coque de petróleo} & \encabezadodorado{GLP} & \encabezadodorado{Gasolinas y naftas} & \encabezadodorado{Querosenos} \\
    \midrule
    \endfirsthead

    \multicolumn{6}{l}{\small\textit{Continuación...}} \\
    \toprule
    \rowcolor{gobmxDorado} \encabezadodorado{} & \encabezadodorado{Coque de carbón} & \encabezadodorado{Coque de petróleo} & \encabezadodorado{GLP} & \encabezadodorado{Gasolinas y naftas} & \encabezadodorado{Querosenos} \\
    \midrule
    \endhead

    \midrule
    \multicolumn{6}{r}{\small\textit{Continúa en la siguiente página...}} \\
    \endfoot

    \bottomrule
    \endlastfoot

    Oferta interna bruta & 16.82 & 119.25 & 297.83 & 1,073.94 & 127.03 \\
    \hspace{1.35em} Oferta total & 16.82 & 121.84 & 297.83 & 1,078.42 & 127.03 \\
    \hspace{2.7em} Producción & 0 & 0 & 0 & 0 & 0 \\
    \hspace{2.7em} Importación & 16.82 & 118.73 & 298 & 1,084 & 126.76 \\
    \hspace{2.7em} Variación de inventarios & 0 & 3.11 & -0.17 & -5.58 & 0.27 \\
    \hspace{1.35em} Exportación & 0 & -2.59 & 0 & -4.48 & 0 \\
    \hspace{1.35em} Energía no aprovechada & 0 & 0 & 0 & 0 & 0 \\
    Total transformaciones & 0 & 15.05 & 168.92 & 514.64 & 66.33 \\
    \hspace{1.35em} Coquizadoras y hornos & 0 & 0 & 0 & 0 & 0 \\
    \hspace{1.35em} Refinerías y despuntadoras & 0 & 32.82 & 18.55 & 441.01 & 66.33 \\
    \hspace{1.35em} Plantas de gas y fraccionadoras & 0 & 0 & 150.37 & 73.63 & 0 \\
    \hspace{1.35em} Centrales eléctricas & 0 & -17.77 & 0 & 0 & 0 \\
    \hspace{2.7em} Carboeléctrica & 0 & 0 & 0 & 0 & 0 \\
    \hspace{2.7em} Térmica convencional & 0 & -17.77 & 0 & 0 & 0 \\
    \hspace{2.7em} Combustión interna & 0 & 0 & 0 & 0 & 0 \\
    \hspace{2.7em} Turbogas & 0 & 0 & 0 & 0 & 0 \\
    \hspace{2.7em} Ciclo combinado & 0 & 0 & 0 & 0 & 0 \\
    \hspace{2.7em} Nucleoeléctrica & 0 & 0 & 0 & 0 & 0 \\
    \hspace{2.7em} Cogeneración & 0 & 0 & 0 & 0 & 0 \\
    \hspace{2.7em} Hidroeléctrica & 0 & 0 & 0 & 0 & 0 \\
    \hspace{2.7em} Geotérmica & 0 & 0 & 0 & 0 & 0 \\
    \hspace{2.7em} Eólica & 0 & 0 & 0 & 0 & 0 \\
    \hspace{2.7em} Solar fotovoltáica & 0 & 0 & 0 & 0 & 0 \\
  \end{xltabular}

  \vspace{0.25em}
  {\small\textit{Continuación Tabla. Tabla 6.3. Matriz origen destino, oferta interna bruta y transformaciones de energéticos secundarios, 2024}}
  \vspace{0.15em}

  \begin{xltabular}{\textwidth}{YZZZZZ}
    \toprule
    \rowcolor{gobmxDorado} \encabezadodorado{} & \encabezadodorado{Diesel} & \encabezadodorado{Combustóleo} & \encabezadodorado{Otros energéticos} & \encabezadodorado{Gas seco} & \encabezadodorado{Energía eléctrica} \\
    \midrule
    \endfirsthead

    \multicolumn{6}{l}{\small\textit{Continuación...}} \\
    \toprule
    \rowcolor{gobmxDorado} \encabezadodorado{} & \encabezadodorado{Diesel} & \encabezadodorado{Combustóleo} & \encabezadodorado{Otros energéticos} & \encabezadodorado{Gas seco} & \encabezadodorado{Energía eléctrica} \\
    \midrule
    \endhead

    \midrule
    \multicolumn{6}{r}{\small\textit{Continúa en la siguiente página...}} \\
    \endfoot

    \bottomrule
    \endlastfoot

    Oferta interna bruta & 506.11 & -496.67 & -0.27 & 2,569.96 & -6.94 \\
    \hspace{1.35em} Oferta total & 540.53 & 0.33 & -0.27 & 2,578.96 & 23.51 \\
    \hspace{2.7em} Producción & 0 & 0 & 0 & 0 & 0 \\
    \hspace{2.7em} Importación & 543.51 & 0.39 & 0 & 2,578.47 & 23.51 \\
    \hspace{2.7em} Variación de inventarios & -2.98 & -0.06 & -0.27 & 0.49 & 0 \\
    \hspace{1.35em} Exportación & -34.42 & -497 & 0 & -9 & -30.45 \\
    \hspace{1.35em} Energía no aprovechada & 0 & 0 & 0 & 0 & 0 \\
    Total transformaciones & 311.03 & 527.88 & 88.9 & -961.45 & 1,381.52 \\
    \hspace{1.35em} Coquizadoras y hornos & 0 & 0 & 0 & 0 & 0 \\
    \hspace{1.35em} Refinerías y despuntadoras & 344.46 & 574.97 & 38.88 & 590.63 & 0 \\
    \hspace{1.35em} Plantas de gas y fraccionadoras & 0 & 0 & 50.03 & 634.1 & 0 \\
    \hspace{1.35em} Centrales eléctricas & -33.43 & -47.09 & 0 & -2,186.18 & 1,381.52 \\
    \hspace{2.7em} Carboeléctrica & 0 & 0 & 0 & 0 & 50.44 \\
    \hspace{2.7em} Térmica convencional & 0 & -46.15 & 0 & -273.55 & 122.59 \\
    \hspace{2.7em} Combustión interna & -5.12 & -0.94 & 0 & -49.06 & 16.91 \\
    \hspace{2.7em} Turbogas & -28.31 & 0 & 0 & -99.85 & 57.94 \\
    \hspace{2.7em} Ciclo combinado & 0 & 0 & 0 & -1,606.02 & 802 \\
    \hspace{2.7em} Nucleoeléctrica & 0 & 0 & 0 & 0 & 44.3 \\
    \hspace{2.7em} Cogeneración & 0 & 0 & 0 & -157.7 & 46.83 \\
    \hspace{2.7em} Hidroeléctrica & 0 & 0 & 0 & 0 & 86.52 \\
    \hspace{2.7em} Geotérmica & 0 & 0 & 0 & 0 & 13.43 \\
    \hspace{2.7em} Eólica & 0 & 0 & 0 & 0 & 72.66 \\
    \hspace{2.7em} Solar fotovoltáica & 0 & 0 & 0 & 0 & 67.88 \\
  \end{xltabular}

  \vspace{0.25em}
  {\small\textit{Continuación Tabla. Tabla 6.3. Matriz origen destino, oferta interna bruta y transformaciones de energéticos secundarios, 2024}}
  \vspace{0.15em}

  \begin{xltabular}{\textwidth}{YZZ}
    \toprule
    \rowcolor{gobmxDorado} \encabezadodorado{} & \encabezadodorado{Total energéticos secundarios} & \encabezadodorado{Total energéticos} \\
    \midrule
    \endfirsthead

    \multicolumn{3}{l}{\small\textit{Continuación...}} \\
    \toprule
    \rowcolor{gobmxDorado} \encabezadodorado{} & \encabezadodorado{Total energéticos secundarios} & \encabezadodorado{Total energéticos} \\
    \midrule
    \endhead

    \midrule
    \multicolumn{3}{r}{\small\textit{Continúa en la siguiente página...}} \\
    \endfoot

    \bottomrule
    \endlastfoot

    Oferta interna bruta & 4,207.06 & 9,641.77 \\
    \hspace{1.35em} Oferta total & 4,785 & 12,353.61 \\
    \hspace{2.7em} Producción & 0 & 7,358.78 \\
    \hspace{2.7em} Importación & 4,790.19 & 5,009.54 \\
    \hspace{2.7em} Variación de inventarios & -5.2 & -14.72 \\
    \hspace{1.35em} Exportación & -577.94 & -2,451.03 \\
    \hspace{1.35em} Energía no aprovechada & 0 & -260.8 \\
    Total transformaciones & 2,112.82 & -2,492.31 \\
    \hspace{1.35em} Coquizadoras y hornos & 0 & 0 \\
    \hspace{1.35em} Refinerías y despuntadoras & 2,107.65 & -11.84 \\
    \hspace{1.35em} Plantas de gas y fraccionadoras & 908.12 & -456.44 \\
    \hspace{1.35em} Centrales eléctricas & -902.95 & -2,024.04 \\
    \hspace{2.7em} Carboeléctrica & 50.44 & -97.02 \\
    \hspace{2.7em} Térmica convencional & -214.88 & -271.96 \\
    \hspace{2.7em} Combustión interna & -38.21 & -43.04 \\
    \hspace{2.7em} Turbogas & -70.22 & -73.05 \\
    \hspace{2.7em} Ciclo combinado & -804.02 & -804.02 \\
    \hspace{2.7em} Nucleoeléctrica & 44.3 & -85.2 \\
    \hspace{2.7em} Cogeneración & -110.87 & -110.87 \\
    \hspace{2.7em} Hidroeléctrica & 86.52 & -175.66 \\
    \hspace{2.7em} Geotérmica & 13.43 & -77.89 \\
    \hspace{2.7em} Eólica & 72.66 & -147.52 \\
    \hspace{2.7em} Solar fotovoltáica & 67.88 & -137.82 \\
  \end{xltabular}
\end{tabladoradoLargo}
\vspace{-4pt}
\fuente{Elaboración SENER}



\section*{Glosario}
\phantomsection
\addcontentsline{toc}{section}{Glosario}

\entradaGlosario{Bagazo de caña}{Energético primario compuesto de material fibroso resultante después de la extracción de jugo en la última molienda, así como otros materiales orgánicos provenientes de la cosecha de la propia caña.}
\entradaGlosario{Biocombustible}{Son los combustibles en estado gaseoso, líquido o sólido que se obtienen a partir del aprovechamiento energético directo de la biomasa o mediante su transformación física, química o biológica.
No se consideran biocombustibles las mezclas con derivados del petróleo (petrolíferos), incluso si contienen una fracción de origen biológico.
}
\entradaGlosario{Biogás}{Es un energético primario en estado gaseoso, producido a partir de la descomposición anaeróbica de biomasa y/o de la fracción biodegradable de los residuos. Su composición permite que, tras un proceso de purificación, pueda usarse como biocombustible o como sustituto del gas natural en aplicaciones energéticas principalmente en generación de energía eléctrica y calefacción.}
\entradaGlosario{Carboeléctrica}{Instalaciones que queman carbón para producir vapor con el fin de generar energía eléctrica.
}
\entradaGlosario{Carbón mineral}{Es un energético primario en estado sólido, de color negro o marrón, que contiene esencialmente carbono, pequeñas cantidades de hidrógeno, oxígeno, nitrógeno, azufre y otros elementos, el cual proviene de la degradación de organismos vegetales durante un largo periodo de tiempo.}
\entradaGlosario{Central Eléctrica}{Instalación destinada a la generación de energía eléctrica a través de la conversión de energía mecánica, obtenida a partir de diversas fuentes de energía primaria, tales como combustibles fósiles, energía hidráulica, nuclear, eólica, solar fotovoltaica, entre otras. Estas instalaciones integran sistemas de conversión energética que incluyen turbinas, generadores eléctricos y equipos auxiliares para transformar y optimizar la producción eléctrica conforme a los requerimientos de la red de suministro.
}
\entradaGlosario{Centros de Transformación}{Instalaciones industriales de transformación de la energía primaria en productos energéticos secundarios con características que permiten su uso o consumo final.
}
\entradaGlosario{Ciclo Combinado}{Tecnología de generación eléctrica que integra dos ciclos termodinámicos complementarios para maximizar la eficiencia energética. Una turbina de gas convierte la energía química del combustible en energía mecánica a alta temperatura y presión, generando electricidad de manera directa. Los gases de escape calientes resultantes de esta turbina no se desperdician, sino que se canalizan hacia una caldera de recuperación de calor, donde se produce vapor de alta presión. Este vapor alimenta una turbina de vapor de condensación, que convierte la energía térmica residual en energía mecánica adicional, incrementando la producción eléctrica total sin aumentar significativamente el consumo de combustible.
}
\entradaGlosario{Cogeneración}{Generación de energía eléctrica producida de manera conjunta con vapor u otro tipo de energía térmica secundaria útil o ambos; producción directa o indirecta de energía eléctrica mediante la energía térmica no aprovechada en los procesos, o generación directa o indirecta de energía eléctrica cuando se utilicen combustibles producidos en los procesos industriales.}
\entradaGlosario{Cogeneración Eficiente}{Es la generación de energía eléctrica usando en conjunto una energía térmica secundaria.}
\entradaGlosario{Combustión interna}{Principio de fundamento empleado en centrales eléctricas que utilizan motores de combustión interna, generalmente basados en la tecnología de motores diesel, y alimentados por combustibles líquidos como combustóleo o diesel. En estos sistemas, la energía química del combustible se convierte en energía térmica mediante un proceso de combustión controlada dentro de los cilindros del motor. La rápida expansión de los gases de combustión generados produce un movimiento lineal o rotativo en los pistones, transformando la energía térmica en energía mecánica. Esta energía mecánica es transmitida a un generador eléctrico que, a través de un sistema electromagnético, convierte el movimiento rotatorio en energía eléctrica.}
\entradaGlosario{Combustóleo}{Energético secundario compuesto por una mezcla compleja de hidrocarburos pesados con una concentración de fracciones residuales de alto peso molecular, generado durante las etapas finales de la refinación del petróleo. Debido a su alta viscosidad y contenido de compuestos sulfatados, es comúnmente empleado en calderas industriales y procesos de generación térmica que toleran combustibles de menor calidad.}
\entradaGlosario{Consumo final}{Es la energía y materia prima que funge como combustible destinado a los diferentes sectores de  la economía para su consumo, de este mismo se deriva el consumo final no energético y el consumo final energético.}
\entradaGlosario{Consumo final energético}{Es el aprovechamiento de la energía primaria y secundaria para atender las necesidades energéticas de los sectores de consumo; Agropecuario, industrial, comercial, residencial público y transporte.}
\entradaGlosario{Consumo final no energético}{Es el uso de energéticos primarios y secundarios como insumos en procesos industriales que no tienen como fin la generación de energía, sino la transformación de materias primas para la elaboración de productos no energéticos. Este tipo de consumo es común en sectores como la industria petroquímica, donde se utilizan hidrocarburos para fabricar plásticos, fertilizantes, solventes u otros productos químicos.
}
\entradaGlosario{Consumo propio del sector energético}{Es el consumo de energía primaria y secundaria utilizado internamente por las propias instalaciones del sector energético, incluyendo plantas de extracción, procesamiento, transformación, transporte y distribución. Este consumo no se destina a los usuarios finales, sino que es necesario para operar los sistemas que hacen posible la generación, conversión y entrega de energía.
}
\entradaGlosario{Coque de carbón}{Es un energético secundario, proveniente de la coquización del carbón bituminoso en plantas de coque. Se produce al someter el carbón siderúrgico a un proceso de calentamiento controlado en ausencia de oxígeno. Durante este tratamiento térmico, se eliminan los componentes volátiles del carbón, dando lugar a un material rico en carbono, con alta resistencia mecánica y capacidad calorífica.}
\entradaGlosario{Coque de petróleo}{Es un energético secundario sólido producido en la etapa de coquización dentro del proceso de refinación del petróleo crudo. Se genera un residuo carbonoso con alto contenido de carbono y bajo en hidrógeno, resultado de la descomposición térmica controlada de fracciones pesadas residuales a temperaturas elevadas.}
\entradaGlosario{Coquizadoras y hornos}{Plantas de procesos donde se obtiene coque de petróleo y carbón como resultado de la combustión del carbón mineral y de otros materiales carbonosos.}
\entradaGlosario{Diesel}{Es un energético secundario en estado líquido obtenido principalmente de la destilación del petróleo a una temperatura entre los 200°C y 380°C, este se llega a clasificar dependiendo a su uso, si es para vehículos, maquinaria agrícola o uso específicamente para calderas o que generan calor debido a su alto contenido de parafinas.
}
\entradaGlosario{Diferencia estadística}{En el balance energético, la diferencia estadística representa la discrepancia numérica entre la oferta total energética y su consumo final. Esta diferencia se calcula restando el consumo final de la oferta total.}
\entradaGlosario{Energía eléctrica}{Es la energía producida por el desplazamiento de cargas eléctricas principalmente electrones, a través de un conductor. Este flujo de cargas se produce como consecuencia de una diferencia de potencial eléctrico entre dos puntos, la cual establece un campo eléctrico que induce el movimiento ordenado de dichas cargas. Este desplazamiento permite la realización de trabajo eléctrico.}
\entradaGlosario{Energía eólica}{Es la energía primaria obtenida a partir de la conversión de la energía cinética del viento, producida por el desplazamiento de las corrientes de aire ocasionadas por el calentamiento no uniforme de la superficie terrestre.}
\entradaGlosario{Energía geotérmica}{Energía primaria extraída del calor interno de la tierra,  concentrado en sistemas o yacimientos geotérmicos subterráneos. Este calor se manifiesta en la superficie a través de manantiales, suelos calientes, volcanes de lodo, fumarolas, géiseres y zonas de alteración hidrotermal.}
\entradaGlosario{Energía hidráulica}{Es la energía primaria obtenida a partir de la transformación de la energía potencial y cinética del agua en movimiento, proveniente de la circulación natural de masas líquidas en ríos o arroyos, que fluyen desde zonas elevadas o montañosas hacia áreas de menor latitud, como llanuras o costas. Esta energía es aprovechada para generar electricidad mediante sistemas hidráulicos.}
\entradaGlosario{Energía no aprovechada}{Es la porción de energía disponible en el sistema energético que, debido a restricciones técnicas, económicas o de infraestructura, no puede ser utilizada de manera efectiva. Esta energía puede encontrarse en fuentes que aún no han sido desarrolladas, en excedentes no consumidos o en pérdidas generadas durante los procesos de transformación, transporte o de distribución. Aunque forma parte del balance energético, no contribuye directamente al consumo final.
}
\entradaGlosario{Energía nuclear}{Energía primaria obtenida de los núcleos de los átomos siendo este una partícula más pequeña en la que se puede dividir un material. Un átomo posee dos tipos de partículas, neutrones y protones, los cuales se mantienen unidos por dicha energía. El combustible más utilizado para la obtención de esta energía es el uranio al ser el material más pesado de la naturaleza y con mayor energía atómica.
}
\entradaGlosario{Energía primaria}{Son los energéticos que se obtienen directamente de los recursos naturales sin haber sido sometidos a ningún proceso de transformación. Incluyen fuentes como el petróleo crudo, el gas natural, el carbón, la leña, la energía solar, eólica, hidráulica, entre otras. Esta energía es la base para generar energía secundaria o productos energéticos que pueden ser utilizados directamente en el consumo.
}
\entradaGlosario{Energía secundaria}{Son los energéticos que se obtienen a partir de la transformación de las fuentes de energía primaria en centros especializados. Estos productos energéticos tienen características específicas que hacen aptos para su uso directo en los distintos sectores de consumo. }
\entradaGlosario{Energía solar}{Energía primaria derivada por la radiación electromagnética emitida por el sol. Puede ser captada y convertida en energía térmica para el calentamiento de fluidos o transformada en energía eléctrica mediante tecnologías fotovoltaicas.}
\entradaGlosario{Energías limpias}{Son las fuentes de energía y procesos de generación de electricidad que no producen emisiones o residuos dentro de los umbrales establecidos por las normativas ambientales vigentes. Ayudan a la reducción del impacto negativo ambiental y contribuyen a la sostenibilidad energética.}
\entradaGlosario{Energías renovables }{Son las energías cuyas fuentes residen en fenómenos de la naturaleza, procesos o materiales susceptibles de ser transformados en energía aprovechable que se regeneran naturalmente o con capacidad de regeneración a escala de tiempo del ser humano. Se consideran fuentes de Energías Renovable; al viento tanto en zonas terrestres como marinas, la radiación solar en todas sus formas, el movimiento del agua en cauces naturales o artificiales con embalses ya existentes, la energía oceánica en sus distintas formas, la energía que se obtiene mediante el aprovechamiento del calor interno de la tierra y los energéticos que determine la Ley de Biocombustibles.
}
\entradaGlosario{Eoloeléctrica}{La energía eólica generada se da en las turbinas eólicas o aerogeneradores de eje horizontal o de eje vertical, que son los dispositivos que transforman la energía cinética del viento en energía mecánica para impulsar un generador eléctrico.
}
\entradaGlosario{Exportación}{Corresponde a la venta y envío de energéticos primarios y secundarios a mercados fuera del territorio nacional, destinada para su uso o consumo en otros países. Forma parte del balance energético al representar la energía saliente del país, reduciendo la cantidad disponible para su consumo interno.
}
\entradaGlosario{Gas Licuado de Petróleo (GLP)}{Es un energético secundario obtenido de los procesos de refinación del petróleo y de las plantas procesadoras de gas natural. Consiste en una mezcla de hidrocarburos gaseosos, principalmente propano y butano, que se encuentran disueltos en el petróleo o forman parte del gas natural, y que se almacenan y transportan en estado líquido bajo presión. 
}
\entradaGlosario{Gas Natural}{Energía primaria constituida por una mezcla de gases que se obtienen directamente de la extracción de yacimientos subterráneos o resultado de procesos industriales de separación. Su componente principal es el metano, acompañado de otros hidrocarburos como el etano, propano, butano y pentanos. Así mismo, puede contener dióxido de carbono, nitrógeno y ácido sulfhídrico, entre otros.}
\entradaGlosario{Gas seco}{Energético secundario derivado del tratamiento del gas natural o de la refinación de petróleo, caracterizado por estar libre de hidrocarburos condensables, compuestos principalmente por metano en estado gaseoso.
}
\entradaGlosario{Gasóleo}{Energético secundario constituido por fracciones intermedias del proceso de refinación de petróleo, empleado principalmente como combustible en motores diesel, quemado en sistemas de calefacción central y como carga de alimentación para la industria química.
}
\entradaGlosario{Gasolina}{Energético secundario formado por la mezcla de hidrocarburos líquidos volátiles, entre los que predominan las parafinas ramificadas, aromáticas, naftenos y olefinas. Su principal uso es como combustible en motores de combustión interna, debido a su alta capacidad de liberación de energía durante la ignición. }
\entradaGlosario{Generación Distribuida}{Generación de energía eléctrica a través de generadores exentos, es decir, centrales eléctricas de pequeña escala que no requieren ni cuentan con un permiso de generación otorgado por la autoridad reguladora pero puede estar interconectada a un circuito de distribución. Esta generación se lleva a cabo generalmente cerca del punto de consumo, y está destinada a abastecer parcial o totalmente la demanda local.}
\entradaGlosario{Geotermoeléctrica}{Funciona de una manera similar a una planta termoeléctrica convencional. Consta de cuatro elementos fundamentales: una caldera, para hervir agua y generar vapor con alta presión y temperatura; una turbina, vayas hojas o álabes se mueven al ser impulsados por ese vapor; un generador, que recibe el movimiento de los álabes de la turbina y que convierte la energía mecánica en energía eléctrica y una subestación, cuyo transformador eléctrico eleva el voltaje de la energía producida en el generador, hasta alcanzar la tensión requerida para su transmisión. Estas plantas no requieren de ningún combustible fósil o nuclear, ni requieren calderas para producir vapor de agua, sino que solo aprovechan el producido por los yacimientos geotérmicos.
}
\entradaGlosario{Hidroeléctrica}{La energía hidráulica se transforma en electricidad al circular agua a presión a través de una turbina para producir energía mecánica. Las turbinas transmiten la energía mecánica de su rotación, mediante un eje para finalmente llegar a un generador de electricidad. Estas hidroeléctricas acumulan grandes cantidades de agua por medio de presas. Estas se utilizan para la generación en horas pico, debido a que pueden garantizar cantidades considerables de electricidad de forma constante durante ciertos periodos.}
\entradaGlosario{Importación}{Es la adquisición de energéticos primarios y secundarios provenientes de países extranjeros, los cuales se integran a la oferta total de energía disponible en el país.
}
\entradaGlosario{Infraestructura energética}{Conjunto integral de instalaciones, sistemas y tecnologías de gran escala que habilitan el desarrollo, operación, control y mantenimiento de los sectores estratégicos de energía. Esta infraestructura abarca los flujos energéticos del país, incluyendo tanto la infraestructura del Sector Eléctrico como la del Sector Hidrocarburos.}
\entradaGlosario{Leña}{Materia prima proveniente de vegetación forestal maderable que se utiliza como energético primario. Su uso principal es como combustible , aunque también puede destinarse a la producción de celulosa, tableros y carbón vegetal.}
\entradaGlosario{Naftas}{Es un energético secundario, derivado de la refinación del petróleo crudo, aunque también se extrae por la destilación del proceso de separación de los líquidos del gas húmedo en los Centros Procesadores de Gas y son clasificados en ligeras y pesadas, en función de su temperatura inicial de ebullición.
}
\entradaGlosario{Nucleoeléctrica}{Instalación de producción de energía eléctrica, basada en la energía nuclear y en la fusión nuclear para generar electricidad.}
\entradaGlosario{Oferta interna bruta}{Es la suma de la oferta total, compuesta por la producción nacional de energéticos primarios, las importaciones y la variación de inventarios, menos las exportaciones y la energía no aprovechada . Representa la totalidad de energía disponible en el país para el consumo interno.}
\entradaGlosario{Oferta total}{Es la suma algebraica de la producción de energéticos primarios, las importaciones y la variación de inventarios. Representa la cantidad total de energía disponible antes de contabilizar exportaciones y la energía no aprovechada.}
\entradaGlosario{Pérdidas no técnicas}{Son las pérdidas de energía eléctrica que ocurren cuando parte del consumo no es registrado por los sistemas de medición. Estas pérdidas no están relacionadas con las propiedades físicas del sistema, sino que se originan por causas como el uso ilícito de la energía, fallas o manipulaciones en los equipos de medición, errores administrativos, omisiones en el registro de datos o facturación incorrecta.}
\entradaGlosario{Pérdidas técnicas por transporte, transmisión y distribución por energético}{Son las mermas de energía que ocurren  a lo largo de la cadena de suministro, desde la producción hasta el consumo final. Estas pérdidas se originan por factores físicos inherentes al transporte, transmisión o distribución de los energéticos, como fricción, disipación térmica, fugas o resistencia eléctrica. En el caso de los petrolíferos, estas pérdidas suelen integrarse dentro del consumo propio del sector energético. En cambio, para la energía eléctrica, se clasifican en dos tipos: Pérdidas técnicas y las pérdidas no técnicas.
}
\entradaGlosario{Pérdidas técnicas por transporte, transmisión y distribución por energético}{Son las pérdidas de energía que se producen de forma inevitable durante el proceso de transmisión y distribución de la electricidad. Estas pérdidas se manifiestan principalmente como energía térmica generada por el calentamiento de los conductores, transformadores y demás componentes del Sistema Eléctrico Nacional cuando la corriente eléctrica circula a través de ellos. Este fenómeno se conoce como efecto Joule, y es una consecuencia directa de la resistencia eléctrica de los materiales. Las pérdidas técnicas forman parte del funcionamiento normal del sistema eléctrico. }
\entradaGlosario{Petróleo}{Energético primario no renovable, en su forma líquida se presenta asociado a capas de gas natural, en yacimientos que han estado enterrados durante millones de años o cubiertos por los estratos superiores de la corteza terrestre. Son hidrocarburos pesados, los cuales son compuestos orgánicos que se obtienen de la extracción de la perforación de pozos.}
\entradaGlosario{Petrolíferos}{Son los energéticos secundarios que se obtienen de la refinación del petróleo o del procesamiento del Gas Natural y que derivan directamente de hidrocarburos, tales como gasolinas, diesel, querosenos, combustóleo y GLP, entre otros, distintos de los Petroquímicos. Son los energéticos secundarios que se obtienen de la refinación del petróleo o del procesamiento del Gas Natural y que derivan directamente de hidrocarburos, tales como gasolinas, diesel, querosenos, combustóleo y GLP, entre otros, distintos de los Petroquímicos. }
\entradaGlosario{Petroquímicos}{Son aquellos líquidos o gases que se obtienen del procesamiento del Gas Natural o de la refinación del petróleo y su transformación, que se utilizan habitualmente como materia prima para la industria.}
\entradaGlosario{Plantas de gas y fraccionadoras}{Son los centros procesadores de gas natural ya que se encargan de separar los componentes de este energético primario y de los condensados para obtener los energéticos secundarios como el gas seco, gasolinas y naftas, butano, propano, etano y productos no energéticos.}
\entradaGlosario{Producción}{Es la energía obtenida directamente de los energéticos primarios, que luego es sometida a procesos de transformación para generar bienes energéticos destinados al consumo.
}
\entradaGlosario{Producción primaria}{Proceso mediante el cual los energéticos primarios o materias primas se transforman en productos o energéticos secundarios, listos para su uso o consumo.}
\entradaGlosario{Querosenos}{Energético secundario en estado líquido, compuesto por una fracción de petróleo que se obtiene mediante la destilación en el rango de temperaturas entre 150 y 300 °C. Se utiliza comúnmente como combustible en calefacción, lámparas y en aviación (como base del combustible para turbinas).}
\entradaGlosario{Refinerías y despuntadoras}{Son las plantas de proceso del petróleo, donde se encarga de separar este energético primario en sus diferentes subproductos siendo los energéticos secundarios como el GLP, gasolinas y naftas, querosenos, diesel,  combustóleo y productos no energéticos.}
\entradaGlosario{Sector Agropecuario}{Es clasificado como un sector económico primario, derivado a que se encarga de la obtención de materias primas relacionadas con la agricultura y ganadería, en este sector corresponde el consumo energético relacionado con las actividades de bombeo de agua y riego, combustibles utilizados para la agricultura mecanizada y en la ganadería, entre otros.
}
\entradaGlosario{Sector Comercial}{Se clasifica como sector económico terciario, derivado a que se refiere a todas las actividades relacionadas con la compra y venta de productos terminados, ya que este sector no se enfoca al proceso de manufactura si no que es intermediario entre los productores y consumidores finales. En este sector el consumo energético se enfoca en los locales comerciales, restaurantes, hoteles, bienes y servicios, entre otros.
}
\entradaGlosario{Sector Industrial}{Es el sector económico secundario, derivado a que se encarga de la transformación de materias primas en productos de consumo, en este sector corresponde el consumo energético relacionado con las actividades de manufactura entre otros.
}
\entradaGlosario{Sector Público}{Se le conoce como el conjunto de instituciones y organismos gestionados directa o indirectamente por el estado. Este se encarga de proporcionar servicios públicos y administración de recursos beneficiosos para la sociedad, entre estos servicios está el consumo energético en el alumbrado público, bombeo de agua potable y aguas residuales, entre otros, correspondientes al sector energético. }
\entradaGlosario{Sector Residencial}{Corresponde a las propiedades destinadas principalmente a la vivienda, abarca aspectos técnicos, económicos y de consumo energético relacionados con la vivienda urbana y rural principalmente en cocción de alimentos, calentamiento de agua, calefacción, iluminación, refrigeración, planchado entre otros.}
\entradaGlosario{Sector Transporte}{Su clasificación se adentra al sector económico terciario, derivado a que se refiere a todas las actividades y servicios relacionados con el desplazamiento de personas y mercancías de un punto a otro. En este sector el consumo energético se enfoca en el autotransporte, aéreo, ferroviario, marítimo.}
\entradaGlosario{Sistema Eléctrico Nacional}{Es el conjunto de infraestructura que permite suministrar energía eléctrica a los usuarios finales del país, a través de un sistema integrado por la Red Nacional de Transmisión, la Red Nacional de Distribución, las centrales eléctricas, los equipos e instalaciones del CENACE destinados al control operativo del SEN, así como demás elementos que determine la Secretaría.}
\entradaGlosario{Solar fotovoltaica}{Tecnología que convierte directamente la radiación solar en energía eléctrica mediante el efecto fotoeléctrico. Está compuesta por un conjunto de módulos solares, cada uno formado por múltiples celdas fotovoltaicas fabricadas con materiales de semiconductores, principalmente de silicio monocristalino, policristalino o tecnologías de capa delgada, Estos módulos se interconectan eléctricamente en configuraciones en serie y paralelo para alcanzar los niveles de voltaje y corriente requeridos, produciendo electricidad en forma de corriente continua (CC), la energía generada se canaliza hacia inversores que convierten la corriente continua a corriente alterna (AC) para su integración a la red eléctrica o su uso directo.
}
\entradaGlosario{Térmica convencional}{Es una central de vapor que produce electricidad a través de una fuente de calor.
}
\entradaGlosario{Transformación}{Proceso correspondiente a la modificación física, química y/o bioquímica de un energético primario para convertirlo en un energético secundario. Este proceso se lleva a cabo en plantas o centrales de transformación, donde los recursos naturales son tratados para obtener productos energéticos con características específicas que permiten su uso en los distintos sectores de consumo. }
\entradaGlosario{Turbogas}{Sistema de generación eléctrica que utiliza una turbina de gas como elemento central, en la cual la energía cinética de los gases calientes, resultantes de la combustión controlada de aire y combustible, se aprovecha directamente para producir trabajo mecánico. En este proceso, el aire comprimido se mezcla y quema con el combustible en la cámara de combustión, generando gases a alta temperatura y presión que se expanden a través de los álabes de la turbina. Esta expansión convierte la energía térmica y de presión de los gases de energía cinética como impulso a la rotación del eje de la turbina. El movimiento rotatorio generado se acopla a un generador eléctrico, donde se transforma en energía eléctrica. }
\entradaGlosario{Turbosina}{Energético secundario derivado del destilado intermedio del petróleo. Es un combustible específico para aviación, utilizado principalmente en turbinas de aviones por su alta estabilidad térmica y capacidad de combustión.}
\entradaGlosario{Variación de inventarios}{Diferencia entre el valor de los inventarios iniciales y finales de un periodo contable. Esta variación influye en la gestión financiera, afectando costos energéticos, niveles de producción energética y la capacidad de satisfacer la demanda energética. 
}

\printbibliography

\section*{Siglas y Acrónimos}
\phantomsection
\addcontentsline{toc}{section}{Siglas y Acrónimos}

\entradaSigla{ASA}{Aeropuertos y Servicios Auxiliares}
\entradaSigla{BNE}{Balance Nacional de Energía}
\entradaSigla{CFE}{Comisión Federal de Electricidad}
\entradaSigla{CNE}{Comisión Nacional de Energía}
\entradaSigla{CPEUM}{Constitución Política de los Estados Unidos Mexicanos}
\entradaSigla{CRE}{Comisión Reguladora de Energía}
\entradaSigla{LPTE}{Ley de Planeación y Transición Energética}
\entradaSigla{LSE}{Ley del Sector Eléctrico}
\entradaSigla{LSH}{Ley del Sector Hidrocarburos}
\entradaSigla{PEMEX}{Petróleos Mexicanos}
\entradaSigla{PIB}{Producto Interno Bruto}
\entradaSigla{PND}{Plan Nacional de Desarrollo}
\entradaSigla{PROSENER}{Programa Sectorial de Energía}
\entradaSigla{RGD}{Redes Generales de Distribución}
\entradaSigla{RNT}{Red Nacional de Transmisión}
\entradaSigla{SEN}{Sistema Eléctrico Nacional}
\entradaSigla{SENER}{Secretaría de Energía}
\entradaSigla{SNIEn}{Sistema Nacional de Información Energética}


\end{document}
